\documentclass[a4paper,12pt]{article}

% Pakete
\usepackage[utf8]{inputenc}
\usepackage[T1]{fontenc}
\usepackage[german]{babel}
\usepackage{amsmath, amssymb}
\usepackage{geometry}
\usepackage{tikz}
\usetikzlibrary{decorations.pathmorphing}

\geometry{a4paper, margin=2.5cm}

\title{\Huge \textbf{Ich mache einen anderen Titel}}
\author{}
\date{}

\begin{document}
\maketitle
\thispagestyle{empty}
\newpage

\tableofcontents
\newpage



\section{Mechanik von Massenpunkten}

\textbf{Definition Impuls:}
\[
\vec{p} = m \vec{v}
\quad \text{(Impuls)}
\]
\[
[p] = \mathrm{kg\, m\, s^{-1}} = \mathrm{N\, s}
\]

\vspace{0.5cm}

\textbf{Newton: Hallo ich mache noch etwas anders}
Er hat erkannt, dass $\vec{p}$ zur Beschreibung von Kraftwirkungen zentral ist.

\subsection*{1. Newton}
Ist die resultierende Kraft null:
\[
\vec{F}_{\text{res}} = 0 \Rightarrow \vec{p} = \text{const}
\]

\begin{itemize}
\item $\vec{p}=0 \Rightarrow \vec{v}=0$ (Ruhe)
\item $\vec{p}=\text{const}\neq0 \Rightarrow$ Richtung + Betrag konstant
\end{itemize}

\subsection*{3. Newton}
Kräfte treten paarweise auf:
\[
\vec{F}_{12} = -\vec{F}_{21}
\]

\textbf{Beispiel: Gewichtskraft}
\begin{center}
\vspace*{0.5cm}
\begin{tikzpicture}
\draw (0,0) circle (1);
\draw[->, red, thick] (0,0) -- (0,-1.5) node[below] {$\vec{F}_G$};
\node at (0,0) {$m$};
\end{tikzpicture}
\end{center}
\[
\vec{F}_G = m \vec{g}
\]

\vspace{0.5cm}

\subsection*{Impulserhaltung}

\begin{tikzpicture}
\draw[thick] (0,0) -- (6,0);

\draw (1,0) rectangle (2,1);
\node at (1.5,0.5) {$m_1$};
\draw[->] (2,0.5) -- (3,0.5) node[right] {$v_1$};

\draw (4,0) rectangle (5,1);
\node at (4.5,0.5) {$m_2$};
\draw[->] (4,0.5) -- (3,0.5) node[left] {$v_2$};
\end{tikzpicture}

\[
\vec{p}_1 + \vec{p}_2 = \vec{p}_1' + \vec{p}_2'
\]

\[
\Delta \vec{p}_1 + \Delta \vec{p}_2 = 0
\]

\vspace{0.5cm}

\subsection*{Feder (Hooke)}

\begin{tikzpicture}
\draw (0,2) -- (0,3);
\draw[decorate, decoration={zigzag}] (0,2) -- (0,0);
\draw ( -0.5,0) rectangle (0.5,-0.5);
\node at (0,-0.25) {$m$};
\draw[<->] (1,2) -- (1,0) node[midway,right] {$x$};
\end{tikzpicture}

\[
F = -D x
\quad [D] = \mathrm{N/m}
\]

\vspace{0.5cm}

\subsection*{Reibung}

\begin{tikzpicture}
\draw[thick] (0,0) -- (5,0);
\draw (2,0) rectangle (3,1);
\node at (2.5,0.5) {$m$};

\draw[->, blue] (3,0.5) -- (4,0.5) node[right] {$v$};
\draw[->, red] (2,0.5) -- (1,0.5) node[left] {$F_R$};
\draw[->] (2.5,1) -- (2.5,2) node[above] {$F_N$};
\end{tikzpicture}

\[
F_R = \mu_G F_N
\]
\[
F_H \le \mu_H F_N
\]

\vspace{0.5cm}

\subsection*{Schiefe Ebene}

\begin{tikzpicture}
\draw (0,0) -- (5,0) -- (5,2) -- cycle;
\draw (3,1) rectangle (4,2);
\node at (3.5,1.5) {$m$};

\draw[->] (3.5,1.5) -- (3.5,0.2) node[below] {$mg$};
\draw[->] (3.5,1.5) -- (4.5,2) node[right] {$F_N$};
\end{tikzpicture}

\[
F_{\parallel} = mg \sin\alpha, \quad F_N = mg \cos\alpha
\]

Arbeit:
\[
W = \vec{F}\cdot \vec{s} = Fs\cos\alpha
\]

\vspace{0.5cm}

\subsection*{Energie}

\[
E_{\text{kin}} = \frac{1}{2}mv^2
\]

\[
E_{\text{kin}} = -W_{0 \to v}
\]

\[
2as = v^2 - v_0^2
\]

\vspace{0.5cm}

\subsection*{Vollkommen unelastischer Stoß (1D)}

\begin{tikzpicture}
\draw[thick] (0,0) -- (6,0);

\draw (1,0) rectangle (2,1);
\node at (1.5,0.5) {$m_1$};
\draw[->] (2,0.5) -- (3,0.5) node[right] {$v_1$};

\draw (4,0) rectangle (5,1);
\node at (4.5,0.5) {$m_2$};
\draw[->] (4,0.5) -- (3,0.5) node[left] {$v_2$};
\end{tikzpicture}

\[
m_1 v_1 + m_2 v_2 = (m_1+m_2)v'
\]

\[
v' = \frac{m_1 v_1 + m_2 v_2}{m_1+m_2}
\]

\[
E_{\text{kin, vor}} = \frac{1}{2}m_1 v_1^2 + \frac{1}{2}m_2 v_2^2
\]

\[
E_{\text{kin, nach}} = \frac{1}{2}(m_1+m_2)v'^2
\]

\[
E_{\text{kin, vor}} \neq E_{\text{kin, nach}}
\]

\end{document}